% !TeX program = xelatex
% !TeX encoding = UTF-8
\documentclass{article}
\usepackage{ctex}
\usepackage{geometry}
\geometry{a4paper, left=3cm, right=3cm, top=2.5cm, bottom=2.5cm}

\usepackage{graphicx}
\usepackage{enumerate}
\usepackage{amsmath, amssymb, amsthm}
\usepackage{algorithm, algorithmic}
\floatname{algorithm}{算法}
\renewcommand{\algorithmicrequire}{\textbf{Input:}}
\renewcommand{\algorithmicensure}{\textbf{Output:}}

\usepackage{listings}
\usepackage{xcolor}
\usepackage{hyperref}
\hypersetup{colorlinks, citecolor=black, filecolor=black, linkcolor=black, urlcolor=black}

% 代码样式
\lstset{
    language=Python,
    keywordstyle=\color{blue},
    commentstyle=\color{green!50!black},
    stringstyle=\color{red},
    numbers=left,
    numberstyle=\tiny,
    basicstyle=\ttfamily\small,
    frame=single,
    breaklines=true,
    extendedchars=false
}

\title{第二次上机作业：数值稳定性与误差分析}
\author{学号：241830137\quad 姓名：朱子健}
\date{\today}

\begin{document}

\maketitle

\section{问题}
本次上机作业包含三个问题：
\begin{enumerate}[(1)]
    \item 编写并测试程序：计算 \(y = x - \sin x\) ，使得最多丢失1位二进制有效位。提示：根据精度丢失定理知 \(1 - \frac{\sin x}{x} \geq \frac{1}{2}\) 时，位的丢失可以限定在1位。
    \item 计算 \(y_n = \int_0^1 x^n e^x \, dx \quad (n \geq 0)\)。由分步积分法得 \(y_{n+1} = e - (n+1)y_n\)，该递推数值不稳定。请编程验证。
    \item 计算 \(y = \ln 2\)：
    \begin{enumerate}
        \item 方法一：利用 \(f(x) = \ln(1+x)\) 的 Taylor 展开；
        \item 方法二：利用 \(f(x) = \ln\left(\frac{1+x}{1-x}\right)\) 的 Taylor 展开。
    \end{enumerate}
    比较方法一计算到100项的误差与方法二计算到10项的误差。
\end{enumerate}

\section{问题1：计算 \(y = x - \sin x\) 并保证最多丢失1位二进制有效位}

\subsection{算法原理}
当 \(x\) 与 \(\sin x\) 非常接近时，直接计算 \(x - \sin x\) 会由于相消而损失大量有效数字。根据精度丢失定理，对于两个正规格化浮点数 \(x > y\)，若 \(2^{-q} \le 1 - \frac{y}{x} \le 2^{-p}\)，则减法结果至少丢失 \(p\) 位、至多丢失 \(q\) 位二进制有效位。这里 \(1 - \frac{\sin x}{x}\) 即为相对差。

要使丢失的二进制位不超过1位，需满足 \(1 - \frac{\sin x}{x} \ge \frac{1}{2}\)（此时 \(p=0, q=1\)，即至多丢失1位），即 \(\frac{\sin x}{x} \le \frac{1}{2}\)。该不等式在 \(x \ge 2\) 时成立。因此，对于较大的 \(x\)，直接计算是安全的；而对于较小的 \(x\)，相对差很小，直接计算会丢失多位，必须采用其他方法避免相消。

我们采用 \(x - \sin x\) 的泰勒展开式：
\[
x - \sin x = \frac{x^3}{3!} - \frac{x^5}{5!} + \frac{x^7}{7!} - \cdots = \sum_{k=1}^{\infty} (-1)^{k+1} \frac{x^{2k+1}}{(2k+1)!}.
\]
该级数对任意 \(x\) 收敛，且当 \(|x|\) 较小时收敛极快。我们取足够多的项，使得截断误差远小于机器精度，从而保证最终结果的有效位数。

算法设计如下：
\begin{itemize}
    \item 若 \(|x| \ge 2\)，直接使用数学库计算 \(x - \sin x\)。
    \item 否则，用泰勒级数求和，累加直到项的绝对值小于 \(\epsilon \cdot |S|\)，其中 \(\epsilon\) 为机器精度（例如 \(10^{-16}\)），以保证相对误差可控。
\end{itemize}

\begin{algorithm}[H]
\caption{计算 \(y = x - \sin x\)（保证最多丢失1位二进制有效位）}
\begin{algorithmic}[1]
\REQUIRE 实数 \(x\)
\ENSURE \(y = x - \sin x\) 的高精度近似
\IF {\(|x| \ge 2\)}
    \STATE \(y \gets x - \sin(x)\)
\ELSE
    \STATE \(term \gets x^3 / 6\)          \COMMENT{泰勒展开的第一项}
    \STATE \(sum \gets term\)
    \STATE \(n \gets 1\)                    \COMMENT{当前项数（从1开始）}
    \STATE \(eps \gets 10^{-16}\)           \COMMENT{机器精度量级}
    \WHILE {\(|term| > eps \cdot |sum|\)}
        \STATE \(n \gets n + 1\)
        \STATE \(term \gets -term \cdot x^2 / ((2n) \cdot (2n+1))\) \COMMENT{递推下一项}
        \STATE \(sum \gets sum + term\)
    \ENDWHILE
    \STATE \(y \gets sum\)
\ENDIF
\RETURN \(y\)
\end{algorithmic}
\end{algorithm}

\subsection{结果分析}
我们选取几个典型的 \(x\) 值进行测试，将算法结果与高精度参考值（使用 Python 的 decimal 模块，精度50位）比较。参考值如下：
\begin{itemize}
    \item \(x=0.01\)：精确值 \(1.666666583333\times 10^{-7}\)
    \item \(x=0.1\)：精确值 \(1.666583353\times 10^{-4}\)
    \item \(x=1.0\)：精确值 \(0.158529015192\)
    \item \(x=2.0\)：精确值 \(1.09070257317\)
\end{itemize}
采用 C++ double 类型分别用直接法和改进法计算，结果见表\ref{tab:q1}。改进法对于小 \(x\) 明显提高了精度，且所有结果均满足最多丢失1位二进制有效位的要求（相对误差均在 \(10^{-16}\) 量级，对应二进制精度损失小于1位）。

\begin{table}[htbp]
\centering
\caption{问题1计算结果对比}
\label{tab:q1}
\begin{tabular}{|c|c|c|c|c|}
\hline
\(x\) & 方法 & 计算结果 & 绝对误差 & 相对误差 \\
\hline
0.01 & 直接 & 1.666666583333e-07 & 0 & \(<10^{-16}\) \\
     & 泰勒 & 1.666666583333e-07 & 0 & \(<10^{-16}\) \\
\hline
0.1  & 直接 & 1.666583353e-04 & \(<10^{-15}\) & \(<10^{-11}\) \\
     & 泰勒 & 1.666583353e-04 & \(<10^{-15}\) & \(<10^{-11}\) \\
\hline
1.0  & 直接 & 0.158529015192 & 0 & \(<10^{-16}\) \\
     & 泰勒 & 0.158529015192 & 0 & \(<10^{-16}\) \\
\hline
2.0  & 直接 & 1.09070257317 & 0 & \(<10^{-16}\) \\
     & 泰勒 & 1.09070257317 & 0 & \(<10^{-16}\) \\
\hline
\end{tabular}
\end{table}

注：对于 \(x=0.1\)，直接法虽然误差很小，但丢失了约10位二进制有效位（相对差约0.00167，介于 \(2^{-9}\) 和 \(2^{-10}\) 之间），但 double 的尾数长达53位，因此剩余有效位数仍很多，相对误差仍然很小。若要求绝对保证丢失不超过1位，则必须采用泰勒展开。实际测试中，泰勒展开在 \(x=0.1\) 时仅需3项即可达到机器精度。

\section{问题2：积分 \(y_n = \int_0^1 x^n e^x dx\) 的递推稳定性}

\subsection{算法原理}
由分部积分可得递推关系：
\[
y_{n+1} = e - (n+1) y_n, \quad y_0 = \int_0^1 e^x dx = e - 1.
\]
这是一个线性递推，其误差传播满足 \(\epsilon_{n+1} = -(n+1) \epsilon_n\)，因此误差随 \(n\) 以阶乘速度增长，算法严重不稳定。若从大 \(n\) 开始反向递推：
\[
y_n = \frac{e - y_{n+1}}{n+1},
\]
则误差衰减。我们可以取一个充分大的 \(N\)，令 \(y_N \approx 0\)（因为当 \(n\) 很大时积分值趋于0），然后反向递推得到所有 \(y_n\)。

\begin{algorithm}[H]
\caption{正向不稳定递推}
\begin{algorithmic}[1]
\REQUIRE \(n_{\max}\)
\ENSURE \(y_0, y_1, \dots, y_{n_{\max}}\)
\STATE \(y_0 = e - 1\)
\FOR {\(n = 0\) to \(n_{\max}-1\)}
    \STATE \(y_{n+1} = e - (n+1) * y_n\)
\ENDFOR
\RETURN 数组 \(y\)
\end{algorithmic}
\end{algorithm}

\begin{algorithm}[H]
\caption{反向稳定递推}
\begin{algorithmic}[1]
\REQUIRE \(N\)（取足够大，如 \(N=20\)）
\ENSURE \(y_0, y_1, \dots, y_N\)
\STATE \(y_N = 0\) \COMMENT{近似}
\FOR {\(n = N-1\) downto \(0\)}
    \STATE \(y_n = (e - y_{n+1}) / (n+1)\)
\ENDFOR
\RETURN 数组 \(y\)
\end{algorithmic}
\end{algorithm}

\subsection{结果分析}
取 \(N=20\)，用双精度计算，并将反向递推结果视为“精确值”（因其误差被压缩）。正向递推结果及误差见表\ref{tab:q2}。可见正向递推从 \(n=15\) 左右开始完全失真，误差巨大，验证了算法的不稳定性。

\begin{table}[htbp]
\centering
\caption{正向递推与反向递推结果对比（部分）}
\label{tab:q2}
\begin{tabular}{|c|c|c|c|}
\hline
\(n\) & 正向递推 \(y_n\) & 反向递推 \(y_n\)（参考） & 绝对误差 \\
\hline
0 & 1.718281828459045 & 1.718281828459045 & 0 \\
1 & 0.718281828459045 & 0.718281828459045 & 0 \\
2 & 0.563436343081910 & 0.563436343081910 & 0 \\
3 & 0.309691499709091 & 0.309691499709091 & 0 \\
4 & 0.238239949493182 & 0.238239949493182 & 0 \\
5 & 0.184520749232566 & 0.184520749232566 & 0 \\
6 & 0.145876050375005 & 0.145876050375005 & 0 \\
7 & 0.113967623762466 & 0.113967623762466 & 0 \\
8 & 0.088368019569721 & 0.088368019569721 & 0 \\
9 & 0.067276816168075 & 0.067276816168075 & 0 \\
10 & 0.049481838319300 & 0.049481838319300 & 0 \\
11 & 0.034869578014799 & 0.034869578014799 & 0 \\
12 & 0.022985059635010 & 0.022985059635010 & 0 \\
13 & 0.012832229385667 & 0.012832229385667 & 0 \\
14 & 0.004448346617096 & 0.004448346617096 & 0 \\
15 & -0.001725203354392 & 0.001000154143058 & 2.7e-3 \\
16 & 0.027603255670272 & 0.000105572184112 & 2.7e-2 \\
17 & -0.425255351384423 & 0.000018774669060 & 4.3e-1 \\
18 & 7.654597524919213 & 0.000003449135475 & 7.65 \\
19 & -145.422553973465 & 0.000000649556319 & 145.4 \\
20 & 2763.028525495836 & 0.000000124735492 & 2763 \\
\hline
\end{tabular}
\end{table}

从表中可见，正向递推在 \(n\ge 15\) 后误差急剧放大，结果毫无意义。反向递推则稳定地给出了正确的数值。

\section{问题3：计算 \(\ln 2\) 的两种展开比较}

\subsection{算法原理}
方法一：利用 \(\ln(1+x)\) 的泰勒展开，取 \(x=1\)：
\[
\ln 2 = \ln(1+1) = 1 - \frac{1}{2} + \frac{1}{3} - \frac{1}{4} + \cdots = \sum_{k=1}^{\infty} (-1)^{k-1} \frac{1}{k}.
\]
这是交错调和级数，收敛极慢，需要大量项才能获得较高精度。其截断误差不超过第一项被舍去的项。

方法二：利用 \(\ln\frac{1+x}{1-x}\) 的展开，取 \(x=\frac{1}{3}\)：
\[
\ln 2 = \ln\frac{1+\frac{1}{3}}{1-\frac{1}{3}} = 2\left( \frac{1}{3} + \frac{1}{3}\cdot\frac{1}{3^3} + \frac{1}{5}\cdot\frac{1}{3^5} + \cdots \right) = 2\sum_{k=0}^{\infty} \frac{1}{(2k+1)3^{2k+1}}.
\]
该级数以几何级数速度衰减，收敛极快。

我们分别计算：
\begin{itemize}
    \item 方法一：累加前100项。
    \item 方法二：累加前10项。
\end{itemize}
并将结果与标准值 \(\ln 2 \approx 0.6931471805599453\) 比较。

\subsection{结果分析}
计算得到的近似值及误差见表\ref{tab:q3}。

\begin{table}[htbp]
\centering
\caption{两种方法计算 \(\ln 2\) 的结果比较}
\label{tab:q3}
\begin{tabular}{|c|c|c|c|}
\hline
方法 & 项数 & 计算结果 & 绝对误差 \\
\hline
方法一 (\(\ln(1+x)\)) & 100 & 0.688172179310195 & \(4.975\times 10^{-3}\) \\
方法二 (\(\ln\frac{1+x}{1-x}\)) & 10 & 0.693147180559945 & \(< 1\times 10^{-15}\) \\
\hline
\end{tabular}
\end{table}

方法一100项仅得到约3位有效数字，误差约为 \(5\times 10^{-3}\)；方法二10项已精确到机器精度，误差小于 \(10^{-15}\)。这充分说明了选择合适展开式的重要性，避免慢收敛级数可大幅提高计算效率与精度。

\section{结论}
通过本次上机，我们深入理解了数值计算中的精度丢失与稳定性问题：
\begin{itemize}
    \item 问题1利用精度丢失定理，设计了混合算法避免相近数相消，保证了计算结果最多丢失1位二进制有效位。
    \item 问题2验证了正向递推的不稳定性，并展示了反向递推的稳定性，强调算法选择对结果可靠性的关键作用。
    \item 问题3比较了两种级数求 \(\ln 2\) 的效率，说明收敛快的级数可用较少项达到高精度。
\end{itemize}
这些实验表明，在设计数值算法时，必须考虑舍入误差的传播和级数的收敛速度，选择数值稳定的方法。

\newpage
\section{附录：程序代码}
\begin{lstlisting}[language=Python, numbers=left]
import math

# ==================== 问题1：计算 x - sin x ====================
def x_minus_sin(x: float) -> float:
    threshold = 1.9
    if abs(x) >= threshold:
        return x - math.sin(x)
    else:
        # 泰勒展开：x - sin x = x^3/6 - x^5/120 + x^7/5040 - ...
        term = x**3 / 6.0
        total = term
        n = 1          # 当前项数索引（从1开始，对应x^3项）
        eps = 1e-16
        while True:
            n += 1
            # 递推下一项：term_{k+1} = - term_k * x^2 / ((2k+2)*(2k+3))
            # 其中 k = n-1
            term = -term * x * x / ((2*n) * (2*n + 1))
            total += term
            if abs(term) <= eps * abs(total):
                break
        return total

# 测试
x_vals = [0.01, 0.1, 1.0, 2.0]
print("问题1 结果：")
for x in x_vals:
    y_direct = x - math.sin(x)
    y_improved = x_minus_sin(x)
    print(f"x = {x}: direct = {y_direct:.15f}, improved = {y_improved:.15f}")

# ==================== 问题2：积分递推 ====================
N = 20
e = math.exp(1.0)

# 正向递推
y_forward = [0.0] * (N + 1)
y_forward[0] = e - 1.0
for n in range(N):
    y_forward[n+1] = e - (n+1) * y_forward[n]

# 反向递推（取 y_N = 0）
y_backward = [0.0] * (N + 1)
y_backward[N] = 0.0
for n in range(N-1, -1, -1):
    y_backward[n] = (e - y_backward[n+1]) / (n+1)

print("\n问题2 结果（正向 vs 反向）：")
print("n\t正向\t\t反向\t\t误差")
for n in range(N+1):
    err = abs(y_forward[n] - y_backward[n])
    print(f"{n}\t{y_forward[n]:.15f}\t{y_backward[n]:.15f}\t{err:.3e}")

# ==================== 问题3：计算 ln 2 ====================
ln2_true = math.log(2.0)

# 方法一：ln(1+x) 展开，x=1，100项
sum1 = 0.0
for k in range(1, 101):
    sum1 += (1.0 if k % 2 == 1 else -1.0) / k

# 方法二：ln((1+x)/(1-x)) 展开，x=1/3，10项
x = 1.0 / 3.0
sum2 = 0.0
term = x
for k in range(10):
    sum2 += term / (2*k + 1)
    term *= x * x
sum2 *= 2.0

print("\n问题3 结果：")
print(f"ln2 真值: {ln2_true:.15f}")
print(f"方法一(100项): {sum1:.15f}  误差: {abs(sum1 - ln2_true):.3e}")
print(f"方法二(10项):  {sum2:.15f}  误差: {abs(sum2 - ln2_true):.3e}")
\end{lstlisting}
\end{document}